\documentclass[11pt,letterpaper]{article}

\usepackage[top=1in, bottom=1in, left=1in, right=1in, headheight=14pt]{geometry}
\usepackage{fontspec}
\setmainfont{DejaVu Serif}
\setsansfont{DejaVu Sans}
\setmonofont{DejaVu Sans Mono}

\usepackage{xcolor}
\usepackage{titlesec}
\usepackage{fancyhdr}
\usepackage{enumitem}
\usepackage{hyperref}
\usepackage{booktabs}
\usepackage{tabularx}
\usepackage{longtable}
\usepackage{colortbl}
\usepackage{multirow}
\usepackage{array}
\usepackage{parskip}
\usepackage{tcolorbox}
\usepackage{graphicx}
\usepackage{lastpage}

% Technology color scheme: Steel / Electric / Graphite
\definecolor{primary}{HTML}{263238}
\definecolor{secondary}{HTML}{1565C0}
\definecolor{tertiary}{HTML}{37474F}
\definecolor{accent}{HTML}{0288D1}
\definecolor{lightprimary}{HTML}{ECEFF1}
\definecolor{lightsecondary}{HTML}{E3F2FD}
\definecolor{lighttertiary}{HTML}{EFEBE9}
\definecolor{rowgray}{HTML}{F5F5F5}
\definecolor{bordergray}{HTML}{BDBDBD}
\definecolor{subtlegray}{HTML}{757575}
\definecolor{darktext}{HTML}{212121}
\definecolor{alertred}{HTML}{B71C1C}
\definecolor{successgreen}{HTML}{1B5E20}

\titleformat{\section}
  {\Large\bfseries\sffamily\color{primary}}
  {\thesection}{1em}{}
  [\vspace{-0.5em}\textcolor{bordergray}{\rule{\textwidth}{0.4pt}}]

\titleformat{\subsection}
  {\large\bfseries\sffamily\color{secondary}}
  {\thesubsection}{1em}{}

\titleformat{\subsubsection}
  {\normalsize\bfseries\sffamily\color{tertiary}}
  {\thesubsubsection}{1em}{}

\titlespacing*{\section}{0pt}{1.5em}{0.8em}
\titlespacing*{\subsection}{0pt}{1.2em}{0.5em}
\titlespacing*{\subsubsection}{0pt}{0.8em}{0.3em}

\newcommand{\stageheader}[2]{%
  \noindent\colorbox{#2}{\makebox[\textwidth][l]{%
    \textcolor{white}{\bfseries\sffamily\hspace{6pt}#1\hspace{6pt}}}}%
  \vspace{0.5em}\par%
}

\newtcolorbox{asciibox}{
  colback=rowgray, colframe=bordergray,
  arc=1pt, boxrule=0.5pt,
  left=6pt, right=6pt, top=4pt, bottom=4pt,
  fontupper=\small\ttfamily,
  breakable
}

\newtcolorbox{quotebox}{
  colback=lightprimary, colframe=primary,
  arc=2pt, boxrule=0.5pt,
  left=12pt, right=12pt, top=8pt, bottom=8pt,
  breakable
}

\newcommand{\metafield}[2]{\textbf{\sffamily #1} #2\par}

\newcolumntype{L}[1]{>{\raggedright\arraybackslash}p{#1}}
\newcolumntype{C}[1]{>{\centering\arraybackslash}p{#1}}
\newcommand{\tableheader}[1]{\textbf{\sffamily\textcolor{white}{#1}}}

\pagestyle{fancy}
\fancyhf{}
\fancyhead[L]{\small\sffamily\textcolor{subtlegray}{ComfyUI / Skill-Creator Integration}}
\fancyhead[R]{\small\sffamily\textcolor{subtlegray}{GSD Mission Package}}
\fancyfoot[C]{\small\sffamily\textcolor{subtlegray}{Page \thepage\ of \pageref{LastPage}}}
\renewcommand{\headrulewidth}{0.4pt}
\renewcommand{\footrulewidth}{0pt}

\hypersetup{
  colorlinks=true,
  linkcolor=secondary,
  urlcolor=secondary,
  pdfauthor={GSD Ecosystem},
  pdftitle={ComfyUI Integration -- Mission Package},
  pdfsubject={Vision to Mission Pipeline},
}

\begin{document}

% ─── TITLE PAGE ────────────────────────────────────────────────────────────────
\thispagestyle{empty}
\begin{center}
\vspace*{2.5cm}
{\small\sffamily\textcolor{primary}{\textbf{GSD RESEARCH MISSION PACKAGE}}}

\vspace{0.3cm}
\textcolor{bordergray}{\rule{8cm}{0.4pt}}

\vspace{1.5cm}
{\fontsize{26}{32}\selectfont\bfseries\sffamily\textcolor{primary}{ComfyUI Integration\\[0.3em]with Skill-Creator}}

\vspace{0.8cm}
{\Large\sffamily\textcolor{secondary}{Local Image \& Video Production Infrastructure}}

\vspace{0.5cm}
{\large\sffamily\itshape\textcolor{tertiary}{A Deep Research Study}}

\vspace{2.5cm}
{\sffamily\textcolor{accent}{Vision $\rightarrow$ Research $\rightarrow$ Mission}}\\[0.3em]
{\small\sffamily\textcolor{accent}{Full Pipeline Execution}}

\vspace{2.5cm}
{\small\sffamily\textcolor{subtlegray}{Date: March 27, 2026 \quad|\quad Status: Ready for GSD Orchestrator}}\\[0.3em]
{\small\sffamily\textcolor{subtlegray}{Activation Profile: Squadron (12 roles) \quad|\quad Estimated: 4--6 context windows across 2--3 sessions}}
\end{center}
\newpage

% ─── TABLE OF CONTENTS ─────────────────────────────────────────────────────────
\tableofcontents
\newpage

% ═══════════════════════════════════════════════════════════════════════════════
\stageheader{STAGE 1 --- VISION DOCUMENT}{primary}

\section{ComfyUI / Skill-Creator Integration --- Vision Guide}

\metafield{Date:}{March 27, 2026}
\metafield{Status:}{Initial Vision / Research Complete}
\metafield{Depends On:}{gsd-skill-creator v1.49+ (DACP milestone), GSD Chipset Architecture}
\metafield{Context Window:}{4--6 windows (Squadron profile, parallel image + video tracks)}

\subsection{Vision}

The GSD ecosystem has always operated on the Amiga Principle: extraordinary output through
architectural leverage rather than brute force. ComfyUI embodies this principle at the
inference layer --- a directed-acyclic-graph engine where node connections define data flow,
partial re-execution prevents redundant computation, and a clean HTTP/WebSocket API makes
the entire pipeline accessible from any language or framework. Skill-Creator, already managing
the six-step observe-detect-suggest-manage-autoload-learn loop, is the natural orchestrator
for this kind of modular production infrastructure.

The vision is a local, sovereign image and video production layer --- no cloud fees, no
prompt logging, no content filters imposed from outside --- woven directly into the skill-creator
ecosystem via structured DACP bundles. A skill dispatches a three-part bundle (human intent +
workflow JSON + executable trigger code) to a running ComfyUI instance over its REST API,
receives WebSocket progress events in real time, retrieves the finished artifact, and logs
the generation into the skill registry for pattern detection. The loop closes back through
skill-creator's learning layer: recurring prompt patterns become reusable workflow templates;
hardware-specific optimizations accumulate as skill metadata; model availability is
auto-discovered at session start.

This is the \emph{spaces between} made productive. The meaning is not in the image itself
--- it lives in the connection between the user's intent, the workflow graph, the model's
latent space traversal, and the artifact that emerges. Skill-Creator provides the structure;
ComfyUI provides the traversal; the user provides the intent. Nothing is wasted. Only the
parts of the graph that change between executions are re-evaluated.

\subsection{Problem Statement}

\begin{enumerate}
  \item \textbf{Cloud dependency and privacy exposure.} Commercial image/video generation
  APIs log prompts, impose content filters, change pricing without notice, and remove
  user sovereignty over the creative process. Local production eliminates these risks
  but currently requires manual ComfyUI setup outside the GSD workflow.

  \item \textbf{No structured handshake between agents and inference engines.}
  Skill-Creator agents currently have no formal protocol for dispatching generation
  requests to ComfyUI. Workflows are produced ad hoc, results arrive as unstructured
  files, and no learning loop captures which parameters produced which outcomes.

  \item \textbf{Model landscape complexity.} Image models (FLUX.1 Dev/Schnell, SDXL,
  SD 3.5) and video models (LTX-2.3, Wan 2.2, CogVideoX, HunyuanVideo) have radically
  different VRAM profiles, licensing terms, quantization strategies, and ControlNet
  compatibility matrices. No single skill currently navigates this landscape for the
  orchestrator.

  \item \textbf{Security risks from untrusted custom nodes.} ComfyUI's community
  ecosystem has been compromised by malicious nodes carrying cryptominers and credential
  stealers. A GSD-integrated deployment requires an audited node allowlist, Docker
  sandboxing, and safety-by-structure boundaries consistent with the ecosystem's core
  security philosophy.

  \item \textbf{No video production path in the current skill-creator pipeline.}
  Skill-creator can orchestrate text and code tasks but has no primitives for queueing
  multi-second video generation jobs, tracking frame-by-frame progress over WebSocket,
  or assembling multi-clip sequences with consistent seeds.
\end{enumerate}

\subsection{Core Concept}

\textbf{ComfyUI as a Deterministic Inference Microservice} --- a local HTTP/WebSocket
server that receives workflow JSON bundles over \texttt{POST /prompt}, streams node-level
progress events, and returns artifact filenames via \texttt{/history}. Skill-Creator
wraps this with a DACP-compliant skill that translates agent intents into typed workflow
bundles, tracks job state, and feeds outcomes back into the learning registry.

\subsubsection{System Architecture}

\begin{asciibox}
GSD Skill-Creator (Orchestrator Layer)
  |
  |  DACP Bundle: {intent, workflow_json, trigger_code}
  v
ComfyUI Skill (Paula / I/O layer)
  |
  |-- POST /prompt  --> ComfyUI REST API (port 8188)
  |-- ws://8188/ws  --> WebSocket progress feed
  |-- GET /history  --> Artifact retrieval
  |
  v
ComfyUI Server (Local, GPU-bound)
  |
  |-- Image Pipeline: FLUX.1/SDXL/ControlNet/LoRA
  |-- Video Pipeline: LTX-2.3 / Wan 2.2 / CogVideoX
  |-- Upscale/Post:  RealESRGAN / RIFE / RTX-VSR
  |
  v
/output/*.png | /output/*.mp4
  |
  v
Skill Registry (pattern detection, template accumulation)
\end{asciibox}

\subsubsection{Research Module Architecture}

\begin{asciibox}
Module 01: API Bridge        Module 02: Image Pipeline
  ComfyUI REST/WebSocket        FLUX / SDXL / ControlNet
  Workflow JSON schema           LoRA / VAE / Upscale
  Python client patterns        VRAM tier routing
        |                              |
        +----------+  +----------------+
                   |  |
                   v  v
         Module 04: Skill-Creator Integration
           DACP bundle schema for generation
           Skill file structure + trigger code
           MCP bridge to ComfyUI (planning-bridge pattern)
                   |
        +----------+  +--------------+
        |                            |
        v                            v
Module 03: Video Pipeline   Module 05: Security & Sandboxing
  LTX-2.3 / Wan 2.2            Docker isolation strategy
  CogVideoX / Hunyuan           Node allowlist / trust system
  Audio-sync / RIFE interp      GSD safety-by-structure
\end{asciibox}

\subsection{Research Modules}

\subsubsection{Module 01 --- ComfyUI API Bridge}
Architecture deep-dive: the Python/PyTorch backend, Flask-based HTTP server, litegraph.js
frontend, and headless API mode. Key endpoints: \texttt{/prompt} (workflow submission),
\texttt{/ws} (WebSocket progress), \texttt{/history} (artifact retrieval),
\texttt{/object\_info} (node catalog). Workflow JSON schema: DAG structure, node IDs,
\texttt{class\_type}/\texttt{inputs} anatomy. Python client patterns including both
synchronous polling and WebSocket streaming. ComfyScript and ComfyUI-to-Python
transpilation. Partial re-execution optimization (only changed subgraph nodes re-run).

\subsubsection{Module 02 --- Image Production Pipeline}
Current model landscape (March 2026): FLUX.1 Dev (24GB FP16, 12--16GB FP8, 8GB GGUF Q4),
FLUX.1 Schnell (4-step distilled), SDXL (8--12GB), SD 3.5. ControlNet integration, LoRA
loading, VAE selection, KSampler configuration. VRAM tier routing table for auto-selecting
model variant based on hardware profile. Upscale chain: RealESRGAN $\rightarrow$ ESRGAN
variants. NVIDIA NVFP4/FP8 optimization path (2.5x faster, 60\% less VRAM on RTX 50-series).
Batch generation and parameter sweep patterns.

\subsubsection{Module 03 --- Video Production Pipeline}
Text-to-video and image-to-video model matrix: LTX-2.3 (real-time 4K, audio sync,
Day-0 ComfyUI v0.16 support, 24GB native/16GB quantized), Wan 2.2 (best photorealism,
Apache 2.0, 8GB for 1.3B, 24GB for 14B), CogVideoX (beginner-friendly, 8--12GB),
HunyuanVideo 1.5 (13B, 40\% VRAM reduction). Frame interpolation via RIFE/GIMM-VFI.
Multi-clip assembly with consistent seeds. RTX Video Super Resolution node for 4K upscale.
Audio-conditioned generation via LTX-2.3 mel-spectrogram conditioning.

\subsubsection{Module 04 --- Skill-Creator Integration Architecture}
DACP-compliant skill file structure for ComfyUI dispatch: the three-part bundle
(human intent string + workflow JSON payload + Python trigger code). MCP bridge pattern
reusing the planning-bridge Cloudflare tunnel approach. Session handshake for ComfyUI
readiness detection. Skill registry feedback loop: generation metadata (model, steps,
CFG, seed, VRAM peak) stored per generation and surfaced for pattern detection.
Auto-load rules for model-specific workflow templates. ComfyUI MCP server landscape
(artokun/comfyui-mcp, alecc08/comfyui-mcp, nikolaibibo pattern). Claude Code nodes
(christian-byrne/claude-code-comfyui-nodes) as reference for inverse direction.

\subsubsection{Module 05 --- Security and Sandboxing}
Custom node threat surface: eval/exec RCE vectors, path traversal, supply chain
attacks (2025 cryptominer/credential-stealer incidents). ComfyUI's January 2025 security
update: eval/exec bans, pip-install-at-runtime prohibition, verified author system.
Comfy.org CI system for dependency conflict detection. Docker isolation strategy for
node allowlist enforcement. GSD safety-by-structure principle applied: architectural
OS-level sandboxing rather than permission prompts. Secret management (no credentials
in ComfyUI server environment). Network isolation for local-only deployments.

\subsection{Chipset Configuration}

\begin{asciibox}
# comfyui-integration chipset
agnus:           # Orchestration / DMA / context
  model: opus
  role: synthesis, skill-schema design, security decisions
  context_budget: 80k

denise:          # Display / output rendering
  model: sonnet
  role: workflow JSON assembly, documentation, API integration code
  context_budget: 40k

paula:           # Audio / I/O / anticipatory
  model: sonnet
  role: ComfyUI API bridge code, WebSocket handler, DACP bundle builder
  context_budget: 40k

gary_68000:      # CPU / glue logic
  model: haiku
  role: shared schemas, VRAM tier table, node allowlist scaffolding
  context_budget: 10k
\end{asciibox}

\subsection{Scope Boundaries}

\subsubsection{In Scope (v1.0)}
\begin{itemize}[leftmargin=1.5em]
  \item ComfyUI Python API client skill with WebSocket progress tracking
  \item DACP bundle schema for image and video generation requests
  \item VRAM tier routing table (8GB / 12GB / 16GB / 24GB / 40GB+)
  \item Image workflow templates: FLUX FP8, SDXL standard, SDXL+ControlNet
  \item Video workflow templates: LTX-2.3 T2V, Wan 2.2 T2V, CogVideoX T2V/I2V
  \item Skill registry feedback loop: generation metadata logging
  \item Docker isolation spec for ComfyUI server deployment
  \item Node allowlist: curated set of verified, audited custom nodes
  \item MCP bridge integration spec (reusing planning-bridge pattern)
  \item Model auto-discovery at skill startup via \texttt{GET /object\_info}
\end{itemize}

\subsubsection{Out of Scope}
\begin{itemize}[leftmargin=1.5em]
  \item LoRA training or fine-tuning within GSD (separate mission)
  \item Cloud/remote ComfyUI deployment (local-first only in v1.0)
  \item Audio-only generation (ElevenLabs integration is separate)
  \item 3D asset generation (Hunyuan3D, TripoSR --- separate mission)
  \item CivitAI model marketplace integration
  \item Real-time collaborative generation (multi-user queue)
\end{itemize}

\subsection{Success Criteria}

\begin{enumerate}
  \item A skill file in \texttt{.claude/skills/comfyui/} dispatches a generation request
  to a running ComfyUI instance via \texttt{POST /prompt} and returns the artifact path.

  \item DACP bundle schema is formally specified: JSON schema with \texttt{intent},
  \texttt{workflow\_json}, and \texttt{trigger\_code} fields, validated before dispatch.

  \item VRAM tier routing table correctly selects model variant given GPU VRAM
  declaration in skill config (8GB / 12GB / 16GB / 24GB tiers, all four tested).

  \item Image generation round-trip (prompt $\to$ artifact) completes in under
  90 seconds on a 16GB VRAM GPU using FLUX FP8.

  \item Video generation round-trip (prompt $\to$ MP4) completes using at least
  one of LTX-2.3 or Wan 2.2, with progress events streamed over WebSocket.

  \item All ComfyUI custom nodes in the allowlist pass Docker-isolated security
  audit (no eval/exec patterns, no runtime pip install, no non-standard network calls).

  \item Skill registry accumulates generation metadata (model, steps, CFG, VRAM peak,
  latency) and surfaces recurring prompt patterns after 10 generations.

  \item MCP bridge connection between skill-creator and ComfyUI is demonstrated
  (using planning-bridge Cloudflare tunnel or equivalent local transport).

  \item Model auto-discovery at startup populates available checkpoints and nodes
  via \texttt{GET /object\_info} without manual configuration.

  \item Partial re-execution optimization is verified: submitting the same workflow
  with only the seed changed does not reload the checkpoint model.

  \item Security boundary test: a Docker-isolated ComfyUI instance cannot access
  the host filesystem outside its designated \texttt{/models}, \texttt{/input},
  and \texttt{/output} mounts.

  \item Full integration test: a skill-creator agent receives a natural-language
  production brief, translates it to a DACP bundle, dispatches to ComfyUI, and
  delivers a finished image or video artifact without human intervention.
\end{enumerate}

\subsection{Relationship to Other Documents}

\begin{center}
\rowcolors{2}{rowgray}{white}
\begin{tabularx}{\textwidth}{L{4.5cm}X}
\toprule
\rowcolor{primary}
\tableheader{Document} & \tableheader{Relationship} \\
\midrule
GSD Chipset Architecture Vision & ComfyUI skill maps to Paula (I/O); orchestration to Agnus \\
GSD Staging Layer Vision & ComfyUI artifacts flow through staging for review before publication \\
Planning Bridge (MCP) & Reuses Cloudflare tunnel session handshake for ComfyUI dispatch \\
DACP Milestone (v1.49+) & ComfyUI skill implements DACP three-part bundle as first media use case \\
Fox Infrastructure Group & FoxCompute GPU infrastructure is eventual production host \\
Deep Audio Analyzer Mission & Audio-conditioned video (LTX-2.3) bridges audio and video pipelines \\
GSD BBS Educational Pack & ComfyUI-generated imagery for educational visual assets \\
\bottomrule
\end{tabularx}
\end{center}

\subsection{The Through-Line}

The Amiga 500 rendered graphics through a chipset so elegantly partitioned that Agnus,
Denise, and Paula divided the work by domain --- not by brute force, but by architectural
clarity about what each chip was \emph{for}. ComfyUI is the Amiga chipset of local
generative inference. Its DAG is not a bureaucracy of nodes; it is a topology of
dependencies, each edge a declaration that this output is the input to that operation.
Nothing executes that does not need to. Nothing loads that is already resident.

The skill-creator integration does not bolt ComfyUI onto the side of the GSD ecosystem
--- it threads it into the same philosophy. A DACP bundle is not a command; it is a
statement of intent plus structure plus executable form. The ComfyUI server receives
it not as an instruction but as a complete workflow: \emph{here is what you need, here
is how to connect it, here is the code that starts it.} The WebSocket stream that follows
is the system thinking in real time, reporting its traversal of the latent space. The
artifact that arrives at \texttt{/output} is the resolution of a mathematical journey
through compressed probability --- the same journey, at a different scale, that the
versine and exsecant trace between points on a circle.

The spaces between the nodes are where the image lives.

\newpage

% ═══════════════════════════════════════════════════════════════════════════════
\stageheader{STAGE 2 --- RESEARCH REFERENCE}{secondary}

\section{ComfyUI / Skill-Creator Integration --- Research Reference}

\subsection{How to Use This Document}

This reference compiles findings from cold-start web research (March 2026) across five
modules. Primary sources are official repositories, technical documentation, security
advisories, and developer guides from the ComfyUI ecosystem. All claims are attributed
to identified sources.

\textbf{Key source organizations and repositories:}
\begin{itemize}[leftmargin=1.5em]
  \item \textbf{Comfy-Org / ComfyUI} --- \url{https://github.com/comfy-org/ComfyUI} (GPL-3.0)
  \item \textbf{ComfyUI Official Docs} --- \url{https://docs.comfy.org}
  \item \textbf{Comfy.org Blog} --- Security and release announcements
  \item \textbf{artokun/comfyui-mcp} --- Full MCP server with Claude Code plugin
  \item \textbf{alecc08/comfyui-mcp} --- T2I/I2I/resize MCP server for Claude Desktop
  \item \textbf{christian-byrne/claude-code-comfyui-nodes} --- Claude Code as ComfyUI nodes
  \item \textbf{Snyk Labs} --- Custom node security research (Dec 2024 / updated May 2025)
  \item \textbf{Comfy.org Security Blog} --- January 2025 security update post
  \item \textbf{InsiderLLM} --- Local video generation benchmark guide (Feb 2026)
  \item \textbf{NVIDIA Blogs} --- GDC 2026 RTX/ComfyUI optimization announcements
\end{itemize}

\subsection{Module 01 Reference --- ComfyUI API Bridge}

\subsubsection{Architecture Overview}

ComfyUI implements a Python backend (PyTorch inference engine) paired with a Flask-based
HTTP server and a JavaScript frontend (litegraph.js). The critical architectural property
for skill-creator integration is that the UI is merely a frontend for the same HTTP
endpoints any programmatic client uses. Swapping the UI layer does not touch the inference
engine.

The graph executes lazily: nodes evaluate only when their outputs are needed. Partial
re-execution is a first-class feature: only subgraph nodes that changed between
executions re-run. For skill-creator, this means submitting a workflow with a new seed
or prompt does not reload the checkpoint from disk --- model weights stay resident in VRAM.

\subsubsection{Core API Endpoints}

\begin{center}
\rowcolors{2}{rowgray}{white}
\begin{tabularx}{\textwidth}{L{4cm}L{2.5cm}X}
\toprule
\rowcolor{primary}
\tableheader{Endpoint} & \tableheader{Method} & \tableheader{Purpose} \\
\midrule
\texttt{/prompt} & POST & Submit workflow JSON for execution \\
\texttt{/ws} & WebSocket & Real-time node progress events \\
\texttt{/history/\{prompt\_id\}} & GET & Retrieve output filenames + metadata \\
\texttt{/view} & GET & Fetch generated image bytes by filename \\
\texttt{/object\_info} & GET & Full node catalog: types, inputs, defaults \\
\texttt{/free} & POST & Unload cached models, free VRAM \\
\texttt{/queue} & GET/DELETE & Queue status, cancel pending jobs \\
\bottomrule
\end{tabularx}
\end{center}

\subsubsection{Workflow JSON Schema}

A ComfyUI workflow submitted to \texttt{/prompt} is a flat dictionary mapping string
node IDs to node descriptor objects. Each descriptor has a \texttt{class\_type} (the
registered node class name) and an \texttt{inputs} dictionary. Upstream node outputs
are referenced as two-element arrays \texttt{[node\_id, output\_index]}.

\begin{asciibox}
{
  "prompt": {
    "1": {
      "class_type": "CheckpointLoaderSimple",
      "inputs": { "ckpt_name": "flux1-dev-fp8.safetensors" }
    },
    "2": {
      "class_type": "CLIPTextEncode",
      "inputs": { "text": "...", "clip": ["1", 1] }
    },
    "3": {
      "class_type": "KSampler",
      "inputs": {
        "model": ["1", 0],
        "positive": ["2", 0],
        "seed": 42, "steps": 20, "cfg": 3.5,
        "sampler_name": "euler", "scheduler": "normal"
      }
    }
  }
}
\end{asciibox}

\subsubsection{Python Client Patterns}

Two integration patterns are documented in ComfyUI's \texttt{script\_examples}:

\textbf{Pattern A (SaveImage + polling):} Submit workflow $\to$ poll
\texttt{/history/\{id\}} until complete $\to$ fetch image bytes via \texttt{/view}.
Suitable for batch jobs where real-time feedback is not needed.

\textbf{Pattern B (SaveImageWebsocket + streaming):} Submit workflow using
\texttt{SaveImageWebsocket} node $\to$ receive binary image frames directly over
the WebSocket connection. More efficient for real-time applications; no disk I/O
for intermediate results.

\textbf{ComfyScript} (Chaoses-Ib/ComfyScript) provides a Python DSL that transpiles
to workflow JSON, enabling loops, conditionals, and library calls mixed with ComfyUI
node chains --- directly relevant for skill-creator's batch generation and parameter
sweep use cases.

\textbf{ComfyUI-to-Python} (pydn) translates saved workflow JSON files into executable
Python scripts. CLI option \texttt{--output -} pipes the generated image to stdout,
enabling shell-level integration.

\subsection{Module 02 Reference --- Image Production Pipeline}

\subsubsection{Model Landscape (March 2026)}

\begin{center}
\rowcolors{2}{rowgray}{white}
\begin{tabularx}{\textwidth}{L{3.2cm}L{2.2cm}L{2cm}X}
\toprule
\rowcolor{primary}
\tableheader{Model} & \tableheader{VRAM (FP16)} & \tableheader{License} & \tableheader{Notes} \\
\midrule
FLUX.1 Dev & 24 GB & Non-commercial & Best prompt adherence; T5-XXL text encoder (10GB alone) \\
FLUX.1 Dev FP8 & 12--16 GB & Non-commercial & Minimal quality loss vs FP16; recommended for 16GB tier \\
FLUX.1 Dev GGUF Q4 & 8 GB & Non-commercial & Quantized via ComfyUI-GGUF; viable on RTX 3060 12GB \\
FLUX.1 Schnell & 12 GB & Apache 2.0 & 4-step distilled; 5x faster than Dev; commercial use OK \\
FLUX.2 Klein 4B/9B & 8--12 GB & TBD & NVFP4 support; Day-0 ComfyUI; GDC 2026 announcement \\
SDXL Base 1.0 & 8 GB & RAIL-M & Workhorse; thousands of community fine-tunes on CivitAI \\
SD 3.5 Medium & 8 GB & Stability AI & Better text rendering than SDXL; triple CLIP encoding \\
\bottomrule
\end{tabularx}
\end{center}

\subsubsection{VRAM Tier Routing}

Skill-Creator's VRAM tier routing selects the appropriate model variant at startup.
Community benchmarks and official documentation establish these tiers:

\begin{center}
\rowcolors{2}{rowgray}{white}
\begin{tabularx}{\textwidth}{L{2.5cm}L{3cm}X}
\toprule
\rowcolor{primary}
\tableheader{VRAM} & \tableheader{GPU Examples} & \tableheader{Recommended Image Config} \\
\midrule
8 GB & RTX 3060, 4060 & SDXL standard; FLUX GGUF Q4 with \texttt{--lowvram} \\
12 GB & RTX 3060 12GB & SDXL + ControlNet; FLUX FP8 at 768px \\
16 GB & RTX 4080, 4060 Ti 16GB & FLUX.1 Dev FP8 at 1024px (recommended sweet spot) \\
24 GB & RTX 4090, RTX 3090 & FLUX.1 Dev FP16 full precision; SDXL batches \\
\bottomrule
\end{tabularx}
\end{center}

\subsubsection{ControlNet and LoRA Integration}

ControlNet nodes extract structural features (edges, depth, pose) from reference images
and condition the generation process. The ComfyUI node chain: preprocessor node
(e.g., DWPose for pose estimation) $\to$ ControlNet Apply node $\to$ KSampler.
Strength parameter (0--1) tunes influence. LoRA files load via a dedicated
\texttt{LoRALoader} node inserted between \texttt{CheckpointLoaderSimple} and
\texttt{CLIPTextEncode}. Multiple LoRAs stack.

\subsubsection{NVIDIA RTX Optimization Path (March 2026)}

NVIDIA announced at GDC 2026 that NVFP4 and FP8 model variants are now available for
FLUX.2 Klein and LTX-2.3 in ComfyUI, delivering up to 2.5x performance improvement
and 60\% VRAM reduction on RTX 50-series GPUs. RTX Video Super Resolution is available
as a standalone ComfyUI node (4K upscale 30x faster than ESRGAN alternatives). For
the GSD integration, this optimization path should be auto-detected via GPU model
string and applied automatically by the VRAM router skill.

\subsection{Module 03 Reference --- Video Production Pipeline}

\subsubsection{Video Model Matrix (March 2026)}

\begin{center}
\rowcolors{2}{rowgray}{white}
\begin{tabularx}{\textwidth}{L{3cm}L{2.5cm}L{1.8cm}X}
\toprule
\rowcolor{primary}
\tableheader{Model} & \tableheader{Min VRAM} & \tableheader{License} & \tableheader{Strengths} \\
\midrule
LTX-2.3 & 16 GB (24GB full) & Custom & Real-time 4K; audio sync; Day-0 ComfyUI v0.16; best temporal consistency \\
Wan 2.2 & 8 GB (1.3B) & Apache 2.0 & Best photorealism for human subjects; MoE architecture; commercial OK \\
CogVideoX-5B & 8--12 GB & Apache 2.0 & Best for beginners; strong text comprehension; GGUF quantization \\
HunyuanVideo 1.5 & 24 GB & Custom & 13B parameter; 40\% VRAM reduction vs v1.0; strong faces \\
AnimateDiff & 8 GB & Apache 2.0 & Motion modules for SD 1.5/SDXL; access to entire SD ecosystem \\
Mochi 1 & 20 GB & Apache 2.0 & Historical interest; surpassed by Wan 2.2 \\
\bottomrule
\end{tabularx}
\end{center}

\subsubsection{Video Pipeline Patterns}

Three production patterns are established in the community as of early 2026:

\textbf{Rapid Iteration Pipeline (LTX-2.3):} Generate at 768$\times$512 in seconds,
iterate on prompts, then upscale winners with RTX Video Super Resolution 4K node.
Lowest cost-per-iteration path.

\textbf{Quality Pipeline (Wan 2.2 + post-processing):} Generate at 480p/720p,
upscale with RealESRGAN 4$\times$, interpolate frames with RIFE or GIMM-VFI.
Best photorealism; higher per-clip compute.

\textbf{Image-to-Video Pipeline (FLUX/SDXL + CogVideoX I2V):} Generate a hero
image with FLUX or SDXL, animate it with CogVideoX's image-to-video mode.
Strong compositional control; consistent style across clips.

\subsubsection{Audio-Conditioned Video (LTX-2.3)}

LTX-2.3 supports audio-driven video generation by conditioning the diffusion process
on mel-spectrogram features extracted from input audio. ComfyUI v0.16 provides Day-0
support with dedicated LTX 2.3 nodes. Minimum 24GB VRAM native; 16GB with quantized
models at lower resolutions. Workflow supports T2V, I2V, and audio-driven modes
within the same node graph.

For GSD's Deep Audio Analyzer mission, this creates a direct bridge: audio analysis
outputs (tempo, energy envelope, frequency bands) can parameterize the mel-spectrogram
conditioning, producing video that responds to analyzed musical structure.

\subsubsection{Multi-Clip Assembly}

Current models cap at 5--10 second clips. Longer sequences require generating
multiple clips with consistent seeds and overlapping prompts, then assembling via
a post-processing node chain. RIFE frame interpolation smooths transitions.
Quality degrades with naive stitching; keyframing techniques with seed control
maintain consistency.

\subsection{Module 04 Reference --- Skill-Creator Integration Architecture}

\subsubsection{MCP Server Landscape}

A dedicated ecosystem of ComfyUI MCP servers has emerged, all connecting Claude
agents to local ComfyUI instances over stdio/SSE transports:

\begin{center}
\rowcolors{2}{rowgray}{white}
\begin{tabularx}{\textwidth}{L{4cm}X}
\toprule
\rowcolor{primary}
\tableheader{Repository} & \tableheader{Capabilities} \\
\midrule
artokun/comfyui-mcp & Full-featured: workflow execution, visualization, parameter sweep, progress monitor, VRAM watchdog, debug, model management; Claude Code plugin with slash commands \\
alecc08/comfyui-mcp & T2I, img2img, resize; intelligent parameter injection; request history; REST API for MCP tools \\
nikolaibibo/claude-comfyui-mcp & Template-based simple generation; \texttt{comfy\_generate\_simple}, \texttt{comfy\_submit\_workflow}, \texttt{comfy\_get\_status} tools \\
christian-byrne/claude-code-comfyui-nodes & Inverse direction: Claude Code agents as ComfyUI nodes; stateless command execution; context chaining \\
\bottomrule
\end{tabularx}
\end{center}

\subsubsection{DACP Bundle Schema for Generation}

Applying DACP to ComfyUI dispatch produces a three-part bundle:

\begin{asciibox}
DACP Generation Bundle
{
  "intent": {
    "type": "image_generation",
    "description": "Generate hero image for Seattle hip-hop module cover",
    "output_path": ".planning/staging/assets/module05_cover.png"
  },
  "data": {
    "model_tier": "16gb",
    "workflow_template": "flux_fp8_standard",
    "positive_prompt": "...",
    "negative_prompt": "...",
    "width": 1024, "height": 1024,
    "steps": 20, "cfg": 3.5, "seed": -1
  },
  "executable": {
    "skill": "comfyui/dispatch.py",
    "entry": "generate_image",
    "args": "${data}",
    "on_complete": "registry.log_generation"
  }
}
\end{asciibox}

\subsubsection{Planning-Bridge Integration Pattern}

The planning-bridge MCP server (\texttt{deposit\_mission} tool connecting claude.ai to
\texttt{.planning/staging/inbox/}) established the session handshake pattern. The same
Cloudflare tunnel approach can bridge skill-creator's claude.ai session to a local
ComfyUI instance, treating ComfyUI as another planning-layer service. The ComfyUI MCP
server runs as a local process; the Cloudflare tunnel exposes it to the claude.ai
session; the DACP bundle is deposited as a structured JSON file that the MCP skill
picks up and dispatches to \texttt{localhost:8188}.

\subsubsection{Skill Registry Feedback Loop}

Generation metadata logged per job: model name, workflow template, positive/negative
prompt hashes, steps, CFG, seed, VRAM peak (from \texttt{/free} before/after stats),
generation latency, output artifact path. After 10 generations, skill-creator's
pattern detection layer can identify recurring prompt structures and propose workflow
template promotions --- elevating ad-hoc generation parameters to named, reusable
skill entries.

\subsection{Module 05 Reference --- Security and Sandboxing}

\subsubsection{Threat Surface}

ComfyUI custom nodes run as arbitrary Python code with full system access. Snyk Labs
research (December 2024, updated May 2025) identified three primary attack vectors:
(1) \texttt{eval}/\texttt{exec} injection via user-controlled STRING inputs; (2) path
traversal via unvalidated file inputs enabling arbitrary file write; (3) supply chain
attacks --- malicious nodes in the community repository, including cryptominer and
credential-stealer incidents discovered in early 2025.

\subsubsection{ComfyUI Security Roadmap (January 2025 Update)}

The Comfy-Org January 2025 security post committed to: blocking new \texttt{eval}/\texttt{exec}
calls in registry nodes (immediate); prohibiting runtime pip install by custom nodes
(6-month migration window); obfuscation detection (multi-statement patterns, undefined
variables); verified author system; CI-based dependency conflict scanning at
\texttt{ci.comfy.org}; remote kill switch for malicious nodes via ComfyUI Manager;
experimental Windows Sandbox for ComfyUI Desktop. The April 2025 Manager update added
automated scanning for base64 in suspicious contexts, non-standard domain connections,
and unusual filesystem access patterns.

\subsubsection{GSD Sandboxing Specification}

For GSD integration, the security posture follows the ecosystem's
\textbf{safety-by-structure} principle: boundaries are architectural, not advisory.

\begin{center}
\rowcolors{2}{rowgray}{white}
\begin{tabularx}{\textwidth}{L{3.5cm}L{2cm}X}
\toprule
\rowcolor{primary}
\tableheader{Control} & \tableheader{Type} & \tableheader{Implementation} \\
\midrule
Docker isolation & ABSOLUTE & ComfyUI runs in container; host mounts limited to \texttt{/models}, \texttt{/input}, \texttt{/output} \\
Node allowlist & ABSOLUTE & Only nodes passing manual code audit are installed; list version-pinned \\
eval/exec ban & ABSOLUTE & Any custom node using eval/exec is excluded regardless of source reputation \\
No credentials in env & ABSOLUTE & ComfyUI server environment contains no tokens, keys, or credentials \\
Network isolation & GATE & Container has no outbound network except \texttt{huggingface.co} model downloads (during setup only) \\
Runtime pip ban & GATE & No node may install packages at runtime; all deps in \texttt{requirements.txt} \\
\bottomrule
\end{tabularx}
\end{center}

\subsection{Safety and Sensitivity Considerations}

\begin{center}
\rowcolors{2}{rowgray}{white}
\begin{tabularx}{\textwidth}{L{3.5cm}L{2cm}X}
\toprule
\rowcolor{primary}
\tableheader{Consideration} & \tableheader{Type} & \tableheader{Approach} \\
\midrule
Custom node supply chain & ABSOLUTE & Allowlist-only; no untested nodes installed in production \\
Credentials in server env & ABSOLUTE & No tokens or keys in ComfyUI Docker environment \\
Model licensing for commercial use & GATE & FLUX.1 Dev is non-commercial; Schnell/Wan/CogVideoX are Apache 2.0 \\
Content filtering bypass & ANNOTATE & Local deployment has no cloud content filters; user responsibility for acceptable use policy \\
VRAM exhaustion / OOM & ANNOTATE & VRAM watchdog (artokun pattern) logs warnings; clear\_vram before large jobs \\
\bottomrule
\end{tabularx}
\end{center}

\subsection{Source Bibliography}

\subsubsection{Official Repositories and Documentation}
\begin{itemize}[leftmargin=1.5em]
  \item Comfy-Org/ComfyUI. \textit{GitHub Repository}. \url{https://github.com/comfy-org/ComfyUI}
  \item ComfyUI Official Documentation. \url{https://docs.comfy.org}
  \item artokun/comfyui-mcp. \textit{GitHub Repository}. \url{https://github.com/artokun/comfyui-mcp}
  \item alecc08/comfyui-mcp. \textit{GitHub Repository}. \url{https://github.com/alecc08/comfyui-mcp}
  \item christian-byrne/claude-code-comfyui-nodes. \textit{GitHub Repository}. \url{https://github.com/christian-byrne/claude-code-comfyui-nodes}
  \item Chaoses-Ib/ComfyScript. \textit{GitHub Repository}. \url{https://github.com/Chaoses-Ib/ComfyScript}
  \item pydn/ComfyUI-to-Python-Extension. \textit{GitHub Repository}. \url{https://github.com/pydn/ComfyUI-to-Python-Extension}
  \item SaladTechnologies/comfyui-api. \textit{GitHub Repository}. \url{https://github.com/SaladTechnologies/comfyui-api}
\end{itemize}

\subsubsection{Security Research}
\begin{itemize}[leftmargin=1.5em]
  \item Snyk Labs. ``Don't Get Too Comfortable: Hacking ComfyUI Through Custom Nodes.'' December 16, 2024 (updated May 27, 2025). \url{https://labs.snyk.io/resources/hacking-comfyui-through-custom-nodes/}
  \item Yan, Yoland. ``ComfyUI 2025 Jan Security Update.'' Comfy.org Blog, January 3, 2025. \url{https://blog.comfy.org/p/comfyui-2025-jan-security-update}
  \item Doyensec LLC. ``ComfyUI Manager RCE via Custom Node Install.'' Security Advisory, 2025.
  \item Apatero. ``ComfyUI Custom Nodes Security Guide.'' November 26, 2025. \url{https://apatero.com/blog/comfyui-custom-nodes-security-guide-protect-yourself-2025}
\end{itemize}

\subsubsection{Model and Hardware Research}
\begin{itemize}[leftmargin=1.5em]
  \item NVIDIA Blogs. ``NVIDIA and ComfyUI Streamline Local AI Video Generation at GDC.'' March 2026. \url{https://blogs.nvidia.com/blog/rtx-ai-garage-flux-ltx-video-comfyui-gdc/}
  \item InsiderLLM. ``Local AI Video Generation: What Works in 2026.'' February 5, 2026. \url{https://insiderllm.com/guides/local-ai-video-generation/}
  \item AI Free Forever. ``31 Open-Source AI Video Models: Free Tools to Make Videos in 2026.'' \url{https://aifreeforever.com/blog/open-source-ai-video-models-free-tools-to-make-videos}
  \item Apatero. ``Best Open Source Video AI Models 2025: Complete Comparison.'' December 3, 2025. \url{https://apatero.com/blog/open-source-video-models-comparison-december-2025}
  \item Apatero. ``Run AI Image Generator Locally: GPU Guide 2026.'' February 13, 2026. \url{https://apatero.com/blog/run-ai-image-generator-locally-gpu-guide-2026}
  \item Apatero. ``ComfyUI LTX 2.3 Workflow: Audio-Synced Video Tutorial.'' March 2026. \url{https://apatero.com/blog/comfyui-ltx-2-3-audio-video-workflow-tutorial-2026}
  \item Mucahitceylan. ``ComfyUI: A Technical Deep Dive.'' Medium, September 2025. \url{https://medium.com/@mucahitceylan/comfyui-a-technical-deep-dive-into-the-ultimate-stable-diffusion-workflow-engine-df1a7db3f7f5}
  \item Awesome Agents. ``The Best AI Image Generation Models You Can Run on Your Own GPU in 2026.'' February 2026. \url{https://awesomeagents.ai/guides/best-local-image-generation-models-2026/}
\end{itemize}

\subsubsection{Integration and Deployment}
\begin{itemize}[leftmargin=1.5em]
  \item DevLog 20250710. ``ComfyUI API.'' DEV Community. \url{https://dev.to/methodox/devlog-20250710-comfyui-api-1mi0}
  \item BentoML. ``Deploying Custom ComfyUI Workflows as APIs.'' \url{https://www.baseten.co/blog/deploying-custom-comfyui-workflows-as-apis/}
  \item BentoML Documentation. ``ComfyUI: Deploy workflows as APIs.'' \url{https://docs.bentoml.com/en/latest/examples/comfyui.html}
  \item ViewComfy. ``Integrate ComfyUI Workflows into your apps via API.'' April 2025. \url{https://www.viewcomfy.com/blog/integrate-comfyui-workflows-into-your-apps-via-api}
  \item Wong, Shawn. ``How to Use ComfyUI API with Python.'' Medium, March 2025. \url{https://medium.com/@next.trail.tech/how-to-use-comfyui-api-with-python-a-complete-guide-f786da157d37}
\end{itemize}

\newpage

% ═══════════════════════════════════════════════════════════════════════════════
\stageheader{STAGE 3 --- MISSION PACKAGE}{tertiary}

\section{Milestone Specification}

\subsection{Mission Objective}

Design, implement, and validate a DACP-compliant ComfyUI integration skill for the
GSD Skill-Creator ecosystem that enables any Skill-Creator agent to dispatch local
image and video generation requests to a Docker-isolated ComfyUI instance, track
progress over WebSocket, receive artifacts, and accumulate generation metadata in the
skill registry --- all within the Amiga Principle of architectural elegance over
brute force.

\subsection{Architecture Overview}

\begin{asciibox}
Wave 0: Foundation
  [schema] [vram_table] [node_allowlist] [workflow_templates]
                    |
  +-----------------+--------------------+
  |                                      |
Wave 1A: API Bridge          Wave 1B: Production Pipelines
  Module 01 implementation      Modules 02 + 03 implementation
  Python client skill           Image + Video workflow templates
  WebSocket handler             VRAM routing + model selection
  DACP bundle schema            Audio-conditioned video spec
  |                                      |
  +-----------------+--------------------+
                    |
Wave 2: Integration + Security
  Module 04: DACP integration + MCP bridge
  Module 05: Docker spec + node security audit
                    |
Wave 3: Publication + Verification
  Full integration test + documentation
\end{asciibox}

\subsection{System Layers}

\begin{enumerate}
  \item \textbf{Schema Layer (Wave 0):} Shared DACP bundle schema, VRAM tier table, node allowlist,
  workflow template library (JSON files)
  \item \textbf{API Bridge Layer (Wave 1A):} Python skill file, REST client, WebSocket handler,
  artifact retrieval, model auto-discovery
  \item \textbf{Pipeline Layer (Wave 1B):} Image workflow templates (FLUX/SDXL), video workflow
  templates (LTX-2.3/Wan/CogVideoX), VRAM router, upscale chain
  \item \textbf{Integration Layer (Wave 2):} DACP bundle builder, MCP bridge spec, registry
  feedback loop, planning-bridge handshake
  \item \textbf{Security Layer (Wave 2):} Docker Compose spec, node allowlist audit, security
  boundary tests
  \item \textbf{Verification Layer (Wave 3):} Integration test suite, documentation, skill
  registry seeding
\end{enumerate}

\subsection{Deliverables}

\begin{center}
\rowcolors{2}{rowgray}{white}
\begin{tabularx}{\textwidth}{C{0.5cm}L{4cm}XL{2.5cm}}
\toprule
\rowcolor{primary}
\tableheader{\#} & \tableheader{Deliverable} & \tableheader{Acceptance Criterion} & \tableheader{Wave} \\
\midrule
1 & DACP bundle schema & JSON Schema spec passes validation; all three fields required & Wave 0 \\
2 & VRAM tier routing table & Covers 8/12/16/24GB tiers; auto-selects model variant & Wave 0 \\
3 & Node allowlist & $\geq$10 audited nodes; eval/exec and runtime-pip checks complete & Wave 0 \\
4 & Workflow templates & $\geq$5 templates (FLUX FP8, SDXL, SDXL+ControlNet, LTX-2.3 T2V, Wan T2V) & Wave 0 \\
5 & ComfyUI Python skill & Dispatches workflow, tracks WebSocket, retrieves artifact & Wave 1A \\
6 & Model auto-discovery & \texttt{/object\_info} parsed at startup; checkpoints enumerated & Wave 1A \\
7 & Image pipeline & FLUX FP8 round-trip $<$90s on 16GB GPU & Wave 1B \\
8 & Video pipeline & LTX-2.3 or Wan 2.2 T2V completes; MP4 artifact returned & Wave 1B \\
9 & MCP bridge spec & Planning-bridge pattern adapted for ComfyUI; handshake documented & Wave 2 \\
10 & Registry feedback loop & 10-generation test: metadata logged, patterns surfaced & Wave 2 \\
11 & Docker Compose spec & Isolated container: only \texttt{/models}, \texttt{/input}, \texttt{/output} mounts & Wave 2 \\
12 & Security boundary tests & Container cannot access host FS outside mounts (verified) & Wave 2 \\
13 & Integration test & Full brief $\to$ artifact pipeline runs without human intervention & Wave 3 \\
14 & Skill documentation & SKILL.md, README, and workflow template usage guide & Wave 3 \\
\bottomrule
\end{tabularx}
\end{center}

\subsection{Component Breakdown}

\begin{center}
\rowcolors{2}{rowgray}{white}
\begin{tabularx}{\textwidth}{L{3.5cm}L{3cm}L{1.8cm}C{1.5cm}}
\toprule
\rowcolor{primary}
\tableheader{Component} & \tableheader{Dependencies} & \tableheader{Model} & \tableheader{Est. Tokens} \\
\midrule
Shared schema + tables & None & Haiku & 8k \\
DACP bundle schema spec & Shared schema & Sonnet & 12k \\
VRAM tier routing table & Shared schema & Haiku & 6k \\
Node allowlist audit & None & Sonnet & 15k \\
Workflow template library & VRAM table & Sonnet & 20k \\
ComfyUI Python skill (API bridge) & Schema, templates & Sonnet & 25k \\
WebSocket handler + artifact retrieval & Python skill & Sonnet & 15k \\
Model auto-discovery module & API bridge & Sonnet & 10k \\
Image workflow FLUX/SDXL & API bridge, templates & Sonnet & 18k \\
Video workflow LTX/Wan/CogVideoX & API bridge, templates & Sonnet & 20k \\
DACP integration + MCP bridge spec & All Wave 1 components & Opus & 30k \\
Registry feedback loop & MCP bridge, API bridge & Opus & 20k \\
Docker Compose + security spec & Node allowlist & Opus & 18k \\
Security boundary tests & Docker spec & Sonnet & 12k \\
Integration test suite & All components & Opus & 25k \\
Documentation (SKILL.md + README) & All components & Sonnet & 15k \\
\bottomrule
\end{tabularx}
\end{center}

\subsection{Model Assignment Rationale}

\begin{center}
\rowcolors{2}{rowgray}{white}
\begin{tabularx}{\textwidth}{L{2cm}C{1.5cm}X}
\toprule
\rowcolor{primary}
\tableheader{Model} & \tableheader{Share} & \tableheader{Assigned Tasks} \\
\midrule
Opus & $\sim$30\% & DACP integration architecture; MCP bridge design; registry feedback loop; Docker security spec; integration test design; cross-module synthesis \\
Sonnet & $\sim$60\% & Python skill implementation; workflow templates; WebSocket handler; node allowlist audit; image/video pipeline code; documentation \\
Haiku & $\sim$10\% & Shared schema scaffolding; VRAM tier table; JSON fixtures; boilerplate \\
\bottomrule
\end{tabularx}
\end{center}

\section{Wave Execution Plan}

\subsection{Wave Summary}

\begin{center}
\rowcolors{2}{rowgray}{white}
\begin{tabularx}{\textwidth}{C{1cm}L{2.5cm}L{2.5cm}L{1.5cm}X}
\toprule
\rowcolor{primary}
\tableheader{Wave} & \tableheader{Name} & \tableheader{Mode} & \tableheader{Model} & \tableheader{Output} \\
\midrule
0 & Foundation & Sequential & Haiku + Sonnet & Schema, VRAM table, allowlist, workflow templates \\
1A & API Bridge & Parallel & Sonnet & Python skill, WebSocket handler, auto-discovery \\
1B & Pipelines & Parallel & Sonnet & Image + video workflow implementations \\
2 & Integration + Security & Sequential & Opus + Sonnet & DACP integration, MCP bridge, Docker spec \\
3 & Publication & Sequential & Sonnet + Opus & Integration tests, documentation, registry seeding \\
\bottomrule
\end{tabularx}
\end{center}

\subsection{Wave 0: Foundation}

\begin{center}
\rowcolors{2}{rowgray}{white}
\begin{tabularx}{\textwidth}{L{4cm}L{2cm}L{1.8cm}X}
\toprule
\rowcolor{primary}
\tableheader{Task} & \tableheader{Agent} & \tableheader{Model} & \tableheader{Deliverable} \\
\midrule
DACP bundle JSON schema & EXEC-1 & Sonnet & \texttt{schemas/generation\_bundle.json} \\
VRAM tier routing table & EXEC-1 & Haiku & \texttt{config/vram\_tiers.yaml} \\
Node allowlist initialization & EXEC-2 & Sonnet & \texttt{config/node\_allowlist.yaml} (10 audited nodes) \\
Workflow template stubs & EXEC-2 & Haiku & \texttt{templates/*.json} (5 template stubs) \\
\bottomrule
\end{tabularx}
\end{center}

\textbf{HITL Gate:} Review DACP schema and node allowlist before Wave 1 begins. Allowlist
requires Tibsfox approval --- no node installs before manual code audit signoff.

\subsection{Wave 1A: API Bridge (Parallel Track A)}

\begin{center}
\rowcolors{2}{rowgray}{white}
\begin{tabularx}{\textwidth}{L{4cm}L{2cm}L{1.8cm}X}
\toprule
\rowcolor{primary}
\tableheader{Task} & \tableheader{Agent} & \tableheader{Model} & \tableheader{Deliverable} \\
\midrule
ComfyUI Python skill dispatch & EXEC-1 & Sonnet & \texttt{comfyui/dispatch.py} \\
WebSocket progress handler & EXEC-1 & Sonnet & \texttt{comfyui/ws\_monitor.py} \\
Artifact retrieval + history & EXEC-1 & Sonnet & \texttt{comfyui/artifacts.py} \\
Model auto-discovery & EXEC-1 & Sonnet & \texttt{comfyui/discover.py} \\
\bottomrule
\end{tabularx}
\end{center}

\subsection{Wave 1B: Production Pipelines (Parallel Track B)}

\begin{center}
\rowcolors{2}{rowgray}{white}
\begin{tabularx}{\textwidth}{L{4.5cm}L{2cm}L{1.8cm}X}
\toprule
\rowcolor{primary}
\tableheader{Task} & \tableheader{Agent} & \tableheader{Model} & \tableheader{Deliverable} \\
\midrule
FLUX FP8 image workflow & EXEC-2 & Sonnet & \texttt{templates/flux\_fp8.json} (complete) \\
SDXL standard workflow & EXEC-2 & Sonnet & \texttt{templates/sdxl\_standard.json} \\
SDXL + ControlNet workflow & EXEC-2 & Sonnet & \texttt{templates/sdxl\_controlnet.json} \\
LTX-2.3 T2V workflow & EXEC-2 & Sonnet & \texttt{templates/ltx23\_t2v.json} \\
Wan 2.2 T2V workflow & EXEC-2 & Sonnet & \texttt{templates/wan22\_t2v.json} \\
CogVideoX I2V workflow & EXEC-2 & Sonnet & \texttt{templates/cogvideox\_i2v.json} \\
VRAM router implementation & EXEC-2 & Sonnet & \texttt{comfyui/vram\_router.py} \\
\bottomrule
\end{tabularx}
\end{center}

\textbf{HITL Gate:} Review workflow templates before Wave 2. Verify VRAM router selects
correct template given test GPU profiles. Test at least one image generation round-trip.

\subsection{Wave 2: Integration + Security}

\begin{center}
\rowcolors{2}{rowgray}{white}
\begin{tabularx}{\textwidth}{L{4.5cm}L{2cm}L{1.8cm}X}
\toprule
\rowcolor{primary}
\tableheader{Task} & \tableheader{Agent} & \tableheader{Model} & \tableheader{Deliverable} \\
\midrule
DACP bundle builder + validator & CRAFT-API & Opus & \texttt{comfyui/dacp\_builder.py} \\
MCP bridge integration spec & CRAFT-API & Opus & \texttt{docs/mcp\_bridge\_spec.md} \\
Registry feedback loop & CRAFT-API & Opus & \texttt{comfyui/registry\_logger.py} \\
Docker Compose spec & CRAFT-SEC & Opus & \texttt{docker/docker-compose.yml} \\
Security boundary test suite & CRAFT-SEC & Sonnet & \texttt{tests/security/} \\
Node allowlist deep audit & CRAFT-SEC & Opus & \texttt{config/node\_allowlist.yaml} (completed) \\
\bottomrule
\end{tabularx}
\end{center}

\subsection{Wave 3: Publication + Verification}

\begin{center}
\rowcolors{2}{rowgray}{white}
\begin{tabularx}{\textwidth}{L{4.5cm}L{2cm}L{1.8cm}X}
\toprule
\rowcolor{primary}
\tableheader{Task} & \tableheader{Agent} & \tableheader{Model} & \tableheader{Deliverable} \\
\midrule
Full integration test & EXEC-1 & Opus & Test report: brief $\to$ artifact pipeline \\
Core functionality tests & EXEC-2 & Sonnet & \texttt{tests/core/} test suite \\
SKILL.md documentation & EXEC-2 & Sonnet & \texttt{.claude/skills/comfyui/SKILL.md} \\
README and workflow guide & EXEC-2 & Sonnet & \texttt{README.md} \\
Registry seeding (10 test generations) & EXEC-1 & Sonnet & \texttt{.skill\_registry/comfyui/} \\
\bottomrule
\end{tabularx}
\end{center}

\subsection{Cache Optimization Strategy}

\begin{itemize}[leftmargin=1.5em]
  \item Shared DACP schema computed once in Wave 0; injected into all Wave 1/2 contexts as system prompt prefix
  \item VRAM tier table compiled once; distributed to all workflow template tasks as constant reference
  \item Workflow template stubs from Wave 0 serve as cache seeds for Wave 1B completions
  \item Python skill modules (dispatch.py, ws\_monitor.py, artifacts.py) share a common ComfyUI client base class, enabling cache-hit on the base class definition across modules
  \item 5-minute TTL cache optimization: Wave 1A and 1B tasks share the same base prompt prefix (ComfyUI endpoint spec) to maximize prefix cache hits between parallel agents
\end{itemize}

\subsection{Token Budget}

\begin{center}
\rowcolors{2}{rowgray}{white}
\begin{tabularx}{\textwidth}{L{2.5cm}XC{2.5cm}}
\toprule
\rowcolor{primary}
\tableheader{Wave} & \tableheader{Primary Tasks} & \tableheader{Est. Tokens} \\
\midrule
Wave 0 & Schema + tables + templates & 61k \\
Wave 1A & API bridge (3 modules) & 50k \\
Wave 1B & 6 workflow templates + router & 73k \\
Wave 2 & Integration + security (6 modules) & 110k \\
Wave 3 & Tests + documentation & 55k \\
\midrule
\multicolumn{2}{l}{\textbf{Total Estimated}} & \textbf{349k} \\
\bottomrule
\end{tabularx}
\end{center}

\subsection{Dependency Graph}

\begin{asciibox}
[W0: schema] ──────────────────────────────────────┐
[W0: vram_table] ───────────────────────────────────┤
[W0: node_allowlist] ───────────────────────────────┤
[W0: template_stubs] ───────────────────────────────┤
                                                     │
             ┌───────────────────────────────────────┘
             │
  ┌──────────┴──────────┐
  │ [HITL Gate 0]       │
  └──────────┬──────────┘
             │
  ┌──────────┴──────────┐        ┌─────────────────────┐
  │ W1A: API Bridge     │        │ W1B: Pipelines      │
  │  dispatch.py        │        │  flux_fp8.json      │
  │  ws_monitor.py      │        │  wan22_t2v.json     │
  │  artifacts.py       │        │  ltx23_t2v.json     │
  │  discover.py        │        │  vram_router.py     │
  └──────────┬──────────┘        └─────────────────────┘
             │                               │
  ┌──────────┴───────────────────────────────┘
  │ [HITL Gate 1]
  └──────────┬──────────┐
             │
  ┌──────────┴──────────┐
  │ W2: Integration     │
  │  dacp_builder.py    │
  │  mcp_bridge_spec    │
  │  docker-compose.yml │
  │  security tests     │
  └──────────┬──────────┘
             │
  ┌──────────┴──────────┐
  │ W3: Publication     │
  │  integration test   │
  │  SKILL.md           │
  └─────────────────────┘
\end{asciibox}

\section{Test Plan}

\subsection{Test Categories}

\begin{center}
\rowcolors{2}{rowgray}{white}
\begin{tabularx}{\textwidth}{L{3cm}C{1.5cm}L{2cm}L{2cm}X}
\toprule
\rowcolor{primary}
\tableheader{Category} & \tableheader{Count} & \tableheader{Priority} & \tableheader{Failure} & \tableheader{Scope} \\
\midrule
Safety-critical & 6 & Mandatory & BLOCK & Security boundaries, node audit, credentials \\
Core functionality & 18 & Required & BLOCK & API bridge, VRAM routing, artifact retrieval \\
Integration & 10 & Required & BLOCK & DACP bundle, MCP bridge, registry loop \\
Edge cases & 8 & Best-effort & LOG & OOM handling, missing models, malformed workflows \\
\bottomrule
\end{tabularx}
\end{center}

\subsection{Safety-Critical Tests}

\begin{center}
\rowcolors{2}{rowgray}{white}
\begin{tabularx}{\textwidth}{C{1.5cm}L{3.5cm}X}
\toprule
\rowcolor{primary}
\tableheader{ID} & \tableheader{Verifies} & \tableheader{Expected Behavior} \\
\midrule
SC-01 & Node allowlist enforcement & ComfyUI container refuses to load any custom node not on the audited allowlist \\
SC-02 & No eval/exec in allowed nodes & Static scan of all allowlisted node Python files finds zero eval() or exec() calls \\
SC-03 & No runtime pip install & Allowlisted nodes contain no subprocess pip calls or importlib runtime installs \\
SC-04 & Docker filesystem isolation & Container cannot write to or read from host paths outside \texttt{/models}, \texttt{/input}, \texttt{/output} \\
SC-05 & No credentials in environment & \texttt{docker inspect} of ComfyUI container environment contains no API keys, tokens, or passwords \\
SC-06 & No outbound network from container & Container network policy blocks all outbound connections except model download allow-list (setup phase only) \\
\bottomrule
\end{tabularx}
\end{center}

\subsection{Verification Matrix}

\begin{center}
\rowcolors{2}{rowgray}{white}
\begin{tabularx}{\textwidth}{XL{4cm}C{1.8cm}}
\toprule
\rowcolor{primary}
\tableheader{Success Criterion} & \tableheader{Test ID(s)} & \tableheader{Status} \\
\midrule
1. Skill dispatches to ComfyUI via POST /prompt & CF-01, CF-02 & Pending \\
2. DACP bundle schema is formally specified & CF-03, CF-04 & Pending \\
3. VRAM routing covers all four tiers & CF-05, CF-06, CF-07, CF-08 & Pending \\
4. Image round-trip $<$90s on 16GB GPU & CF-09, CF-10 & Pending \\
5. Video round-trip completes with MP4 artifact & CF-11, CF-12 & Pending \\
6. Node allowlist passes Docker security audit & SC-01, SC-02, SC-03 & Pending \\
7. Registry accumulates 10-generation metadata & IN-01, IN-02 & Pending \\
8. MCP bridge connection demonstrated & IN-03, IN-04 & Pending \\
9. Model auto-discovery populates checkpoints & CF-13, CF-14 & Pending \\
10. Partial re-execution optimization verified & CF-15 & Pending \\
11. Docker isolation verified & SC-04, SC-05, SC-06 & Pending \\
12. Full pipeline brief $\to$ artifact & IN-05, IN-06, IN-07 & Pending \\
\bottomrule
\end{tabularx}
\end{center}

\section{Mission Crew Manifest}

\begin{center}
\rowcolors{2}{rowgray}{white}
\begin{tabularx}{\textwidth}{L{2cm}L{3cm}X}
\toprule
\rowcolor{primary}
\tableheader{Role} & \tableheader{Agent} & \tableheader{Responsibility} \\
\midrule
FLIGHT & Orchestrator & Mission direction; go/no-go at each HITL gate \\
PLAN & Planner & Wave decomposition; parallel track optimization; dependency sequencing \\
SCOUT & Research Agent & Additional web research if model landscape changes during execution \\
EXEC-1 & API Bridge Executor & Modules 01 + integration test (dispatch.py, ws\_monitor.py, etc.) \\
EXEC-2 & Pipeline Executor & Modules 02 + 03 (workflow templates, VRAM router) \\
CRAFT-API & API Architect & DACP integration, MCP bridge spec, registry feedback loop \\
CRAFT-SEC & Security Warden & Docker spec, node allowlist audit, security boundary tests \\
VERIFY & Verification Agent & Source validation; schema compliance; test execution \\
INTEG & Integration Agent & Cross-module validation; DACP bundle round-trip testing \\
CAPCOM & HITL Interface & Human approval channel at Wave 0, 1, and 3 gates \\
BUDGET & Resource Manager & Token budget tracking; session boundary planning \\
SURGEON & Health Monitor & Research quality degradation detection; OOM event logging \\
LOG & Chronicle Agent & Audit trail; commit messages; skill registry entries \\
\bottomrule
\end{tabularx}
\end{center}

\section{Execution Summary}

\begin{center}
\rowcolors{2}{rowgray}{white}
\begin{tabularx}{\textwidth}{L{5cm}X}
\toprule
\rowcolor{primary}
\tableheader{Metric} & \tableheader{Value} \\
\midrule
Activation Profile & Squadron (12 roles) \\
Research Modules & 5 (API Bridge, Image Pipeline, Video Pipeline, Integration, Security) \\
Estimated Token Budget & 349k across all waves \\
Estimated Context Windows & 4--6 windows across 2--3 Claude Code sessions \\
Parallel Tracks & 2 (Wave 1A: API Bridge, Wave 1B: Pipelines) \\
HITL Gates & 3 (post-Wave 0, post-Wave 1, post-Wave 3) \\
Safety-Critical Tests & 6 (all BLOCK on failure) \\
Total Test Cases & 42 \\
Primary Deliverable & \texttt{.claude/skills/comfyui/} skill directory \\
Secondary Deliverable & \texttt{docker/} container spec for isolated ComfyUI deployment \\
Licensing Constraint & FLUX.1 Dev = non-commercial; use Schnell/Wan/CogVideoX for commercial work \\
Primary Target GPU & RTX 4090 (24GB) for full pipeline; RTX 4060 Ti 16GB for FLUX FP8 \\
Repository & github.com/Tibsfox/gsd-skill-creator (dev branch) \\
\bottomrule
\end{tabularx}
\end{center}

\vspace{2em}

\begin{quotebox}
\textit{``Each node represents a specific operation, and the connections between them define
the data flow. The beauty of ComfyUI lies in its flexibility --- the graph is not a
command sequence but a topology of dependencies.''}\\[0.5em]
--- Technical summary of ComfyUI's DAG architecture, adapted from community documentation\\[1em]
The GSD skill is not a command either. It is a topology: intent $\rightarrow$ data
$\rightarrow$ executable. The workflow graph is the spaces between the nodes made legible.
When the KSampler traverses the latent space, it is solving the same kind of problem
the versine solves --- reaching across the gap between two known points to find what
lives in the arc between them. The Amiga Principle holds: not more power, but better
connections.
\end{quotebox}

\end{document}
